\documentclass[11pt]{article}
\usepackage[margin=1in]{geometry}
\usepackage{amsmath,amssymb,amsthm,hyperref}
\newtheorem{theorem}{Theorem}
\newtheorem{lemma}[theorem]{Lemma}
\newcommand{\Aut}{\operatorname{Aut}}
\newcommand{\supp}{\operatorname{supp}}
\title{Integrality of the alpha-triangle\\\large MathDB 323219}
\author{Principia Math in Progress}
\date{18 September 2026}
\begin{document}
\maketitle
\noindent\textbf{Status.} Complete mathematical proof. The final published source has
not been available for comparison; campaign eligibility remains pending, with zero
new credit. Review was by the authoring agent only.

\section*{Task 1: Full solution}
Govindarajan's Conjecture 3.1 in \cite{G}, printed page 16, asserts the integrality
of the coefficients defined by
\[
 a_{m+r+1,r}=\frac1{(2m)!!}\sum_{z=0}^{2m}\alpha_{m,z}r^{2m-z},
 \qquad (2m)!!=2^m m!.
\]
Here a Ferrers diagram in $\mathbb Z_{\geq0}^r$ is a finite order ideal for the
coordinatewise order. A strict diagram uses every coordinate. It therefore
contains $\mu_r=\{0,e_1,\ldots,e_r\}$, and $a_{m+r+1,r}$ counts the diagrams
obtained by adding $m$ nodes to $\mu_r$. These interpretations, and the orbit
counting by reduced dimension used below, are prior results in \cite{G},
Sections 2.2--2.5.

\begin{theorem}\label{main}
For every $m\geq0$, the polynomial $2^m m!\,a_{m+r+1,r}$ belongs to $\mathbb Z[r]$.
Consequently every $\alpha_{m,z}$ is an integer. In fact this polynomial is a
nonnegative integer linear combination of the falling factorials
$(r)_k=r(r-1)\cdots(r-k+1)$.
\end{theorem}

For a set $S$ of added nodes, discard the coordinates that are zero on all of
$S$. Its remaining number of coordinates is $k$, the reduced dimension.
Thus $S\subset\mathbb Z_{\geq0}^k$ uses every coordinate, $|S|=m$, every
node of $S$ has coordinate sum at least two, and $\mu_k\cup S$ is a Ferrers
diagram. Let $\Aut(S)\leq\mathfrak S_k$ be its coordinate permutation group.

\begin{lemma}\label{group}
For every such reduced diagram, $|\Aut(S)|$ divides $2^m m!$.
\end{lemma}
\begin{proof}
Let $T\subseteq S$ consist of the nodes with coordinate sum two. Construct a
graph on the $k$ coordinates: $2e_i\in T$ is a loop at $i$, and
$e_i+e_j\in T$ is the undirected edge $\{i,j\}$. Every vertex occurs in this
graph. Indeed, if $v\in S$ uses coordinate $i$, then either $v_i\geq2$, so
$2e_i\leq v$, or $v_i=1$ and some $v_j>0$, so $e_i+e_j\leq v$.
The order-ideal condition supplies the stated quadratic node in each case.
In particular $k\leq2|T|\leq2m$.

The group $G=\Aut(S)$ acts on the $m$ added nodes. Write $K$ for the kernel
and $H$ for the image of this action. Then $|G|=|K||H|$ and $|H|\mid m!$.
Every member of $K$ fixes each quadratic node individually. It therefore fixes
every loop vertex and preserves each edge as an unordered pair.

A vertex on two distinct edges is fixed, as the unique intersection of their
endpoint sets. Along an edge incident with a fixed vertex, the other endpoint
is fixed as well. Hence all vertices of a connected component with a loop,
or with at least three vertices, are fixed. The only components admitting
any remaining motion are components consisting of a single edge and its two
endpoints, with no loops. If there are $q$ such components, restriction to
their endpoints embeds $K$ into $(\mathbb Z/2\mathbb Z)^q$. Notice that
components cannot be permuted by $K$: every edge is fixed individually.
Thus $|K|\mid2^q$; and $q\leq|T|\leq m$. Combining the two divisibilities gives
\[
 |G|=|K||H|\ \mid\ 2^q m!\ \mid\ 2^m m!,
\]
as required. Higher-degree nodes may further restrict $K$, which does not
affect this argument.
\end{proof}

\begin{proof}[Proof of Theorem \ref{main}]
For fixed $m$, let $\mathcal R_m$ contain one representative of each coordinate
permutation orbit of reduced diagrams with $m$ added nodes. It is finite:
$k\leq2m$, and a coordinate value $v_i$ in an added node is at most $m+1$,
since all nodes $2e_i,\ldots,v_i e_i$ must be present.

Fix $S\in\mathcal R_m$ with $k$ active coordinates and $r\geq k$. There are
$\binom rk$ ways to choose its coordinates in $\{1,\ldots,r\}$ and
$k!/|\Aut(S)|$ distinct coordinate labelings on a chosen set. The other
coordinates contribute only their basis nodes in $\mu_r$. Every diagram
counted by $a_{m+r+1,r}$ arises uniquely in this fashion. Therefore
\begin{equation}\label{orbits}
 a_{m+r+1,r}=\sum_{S\in\mathcal R_m}\frac{(r)_{k(S)}}{|\Aut(S)|}.
\end{equation}
This also holds for $0\leq r<k(S)$, since that summand vanishes. It is a
polynomial identity: the right side is a fixed polynomial in $r$, and agrees
with the defining count for every nonnegative integer $r$.
Multiplication by $2^m m!$ and Lemma \ref{group} give
\[
 2^m m!\,a_{m+r+1,r}
 =\sum_{S\in\mathcal R_m}
 \underbrace{\frac{2^m m!}{|\Aut(S)|}}_{\text{nonnegative integer}}(r)_{k(S)}
 \ \in\mathbb Z[r].
\]
Its monomial coefficients are exactly the $\alpha_{m,z}$. For $m=0$, the
empty reduced diagram has $k=0$, trivial automorphism group, and $(r)_0=1$;
the polynomial is $1$, as required. The prescribed zero entries outside the
triangle are integers as well.
\end{proof}

\paragraph{Sharpness and an example.}
The only reduced diagram on $2m$ coordinates is a matching of $m$ disjoint
quadratic edges. To see this, the covering bound $k\leq2|T|\leq2m$ must be
an equality: every added node is quadratic, no loop occurs, and no endpoints
are shared. Its automorphism group has order $2^m m!$. Hence the leading
coefficient of $a_{m+r+1,r}$ is $1/(2^m m!)$, and the normalization is the
smallest positive integer that can make the whole polynomial integral.
For $m=1$, the loop and edge give
\[
 a_{r+2,r}=r+\frac{(r)_2}{2},\qquad
 2a_{r+2,r}=r^2+r.
\]
The same proof applies to any union of coordinate permutation orbits of
reduced diagrams, chosen independently of the ambient dimension. In particular,
it applies to the symmetric-box restrictions in Appendix A of \cite{G}.

\paragraph{Exact verification.}
The accompanying standard-library Python script independently enumerates all
diagrams with $0\leq m\leq4$ in dimensions $1\leq r\leq8$. Their reduced
counts match the source's Table 3. It checks the group, kernel, image, and
orbit identities for all 87 reduced orbits in that range; compares 40
polynomial evaluations with direct counts; and checks 75 additional exact
falling-factorial divisibilities in the source rows $m=5,\ldots,9$.
These finite checks supplement the proof; they are not used to infer the
all-$m$ assertion.

\section*{Task 2: The alpha--beta decomposition}
Investigate whether decomposing reduced diagrams into connected components
proves the finite binomial expansion of $\alpha_{m,z}$ proposed in
Conjecture 3.2 of \cite{G}. The present integrality proof does not establish
that expansion or its asserted vanishing conditions.

\section*{Task 3: Denominators of restricted families}
For families excluding isolated quadratic edges, determine the smallest
normalizing integer. The proof isolates the possible kernel as a group of
endpoint swaps, and may therefore permit sharper bounds than $2^m m!$.

\section*{Task 4: Weighted refinements}
For integer-valued, coordinate-invariant statistics on reduced diagrams,
study the resulting multivariate integral polynomials and identify natural
combinatorial meanings for their falling-factorial coefficients.

\section*{Task 5: Computation by reduced orbits}
Develop enumeration directly on reduced isomorphism types, keeping their
automorphism orders. Formula \eqref{orbits} then recovers counts in every
ambient dimension without enumerating all coordinate labelings.

\begin{thebibliography}{9}
\bibitem{G} S. Govindarajan, \emph{Notes on higher-dimensional partitions},
\href{https://arxiv.org/abs/1203.4419}{arXiv:1203.4419v1} (2012), 36 pages.
Published in \emph{J. Combin. Theory Ser. A} \textbf{120} (2013), 600--622,
\href{https://doi.org/10.1016/j.jcta.2012.11.005}{doi:10.1016/j.jcta.2012.11.005}.
The full preprint was read; the published version has not been read.
\end{thebibliography}
\end{document}
